\section{2026-10-18 — Análisis financiero ajustado a la clase del 03/10}

La clase de modelo de negocio y financiamiento fijó un checklist de doce ítems para el capítulo económico-financiero. Se incorporaron los que faltaban: impuesto a las ganancias del 35\,\% con quebrantos, cargas sociales del empleador (26\,\%), capital de trabajo de dos meses de ingresos recuperado al cierre, tamaño del mercado (TAM, SAM y SOM), sensibilidad de una variable por vez con valor de equilibrio, plan de financiamiento, contraste de la tasa con el CAPM (19,11\,\% con beta igual a 1) y la hoja de supuestos como anexo. Se agregó además un VAN complementario con valor residual.

Por decisión del autor, las horas de desarrollo del año 0 dejan de imputarse al flujo, en línea con el criterio de documento austero de la clase; su valor de mercado queda para una diapositiva de respaldo. La inversión del año 0 baja de 21.133 a 1.741 dólares.

Con esos cambios el escenario neutral pasa de un VAN levemente positivo a uno de $-$12.296 dólares, con una TIR del 14,48\,\%. Se informa tal cual, sin ajustar supuestos para revertirlo; la sensibilidad muestra que el resultado queda en el umbral. Las cifras que el texto cita fuera de las tablas salen ahora de macros generadas por el validador (\texttt{financiero-cifras.tex}), de modo que el capítulo 3 y la conclusión usan los mismos números.

\bigskip
